重组后的核心问题不是资产看起来更多，而是投资者究竟在给哪个法律与会计主体定价。600150是存续吸收方，原601989被吸收；换股、新股登记、上市和会计合并日分别发生在不同时间。当前估值必须使用合并后股本和持续经营范围，历史同比必须使用同一控制下重述比较期。

\section{E02｜法律实施与会计口径时间线}

法律承继、股份登记与会计合并日不能互换。交易于2025年取得审核与注册，原601989于9月5日退市，新股于9月11日登记、9月16日上市；2025年报采用9月30日作为同一控制下企业合并的会计合并日。形式权属登记的逐项完成清单及601989法人注销完成文件，截至截止日未在已归档官方文件中找到。\srcid{CF-10--CF-16, CF-03}

\begin{exhibitbox}[吸收合并关键节点]
\centering
\begin{tikzpicture}[x=1cm,y=1cm,>=Latex,font=\small]
  \draw[navy,very thick,-{Latex}] (0,0) -- (14.7,0);
  \foreach \x/\d/\t in {
    0.5/{2024-09}/{筹划及预案},
    3.0/{2025-07-05}/{上交所审核通过},
    5.6/{2025-07-18}/{证监会注册},
    8.0/{2025-09-05}/{原601989退市},
    10.3/{2025-09-11}/{新股登记},
    12.2/{2025-09-16}/{新增股份上市},
    14.2/{2025-09-30}/{会计合并日}}
  {
    \draw[navy,thick] (\x,-0.12)--(\x,0.12);
    \node[anchor=north,align=center,text width=2.2cm] at (\x,-0.28) {\scriptsize \d\\\t};
  }
\end{tikzpicture}
\sourcenote{CF-06--CF-16；会计合并日CF-03。时间线区分法律、股份与会计节点。}
\end{exhibitbox}

这条时间线的估值含义是：当前市值与前瞻EPS都基于合并后实体；但不能把“实质承继”写成“所有权属登记和法人注销均已完成”。形式手续缺口不否定合并经营范围，却构成治理披露需持续监测的边界。

\section{E03｜股本与EPS分母桥}

合并前600150股本44.7243亿股，新增股份30.5319亿股，合并后总股本75.2562亿股。该股本用于2026E以后EPS和当前市值；2025A基本EPS使用公司披露的加权平均63.2964亿股，不能用期末股本反推替换。

\begin{exhibitbox}[股本与每股口径桥]
\small
\begin{tabularx}{\exhibitboxwidth}{L{4.1cm}R{3.0cm}L{4.4cm}X}
\toprule
\textbf{项目} & \textbf{股数} & \textbf{用途} & \textbf{控制} \\
\midrule
原600150存量股份 & 44.72428758亿股 & 换股前存续股份 & 不作为当前完整EPS分母 \\
向原601989股东新增 & 30.53192530亿股 & 吸收合并换股 & 与存量相加，不重复 \\
合并后期末总股本 & 75.25621288亿股 & 当前市值、2026E--2028E EPS & 唯一前瞻分母 \\
2025A加权平均分母 & 63.29638855亿股 & 公司披露2025A基本EPS & 历史法定EPS保留原分母 \\
\bottomrule
\end{tabularx}
\sourcenote{CF-03、CF-14、CF-15、CF-04。}
\end{exhibitbox}

按33.02元和75.2562亿股计算，当前市值为2,484.96亿元。这个机械桥确保价格、股本、市值与前瞻EPS处于同一A股/CNY口径，也排除了合并前券商报告使用44.7243亿股所产生的虚高EPS。

\section{法定重述、交易备考与旧口径}

2024年原600150收入785.84亿元、归母36.14亿元，是合并前legacy口径；2025年报给出的2024重述比较期收入1,333.51亿元、归母42.20亿元，才是2025同比的正确基数。2025收入1,519.78亿元相对重述2024增长13.97\%，并非相对旧口径看起来的约93\%。\srcid{CF-02, CF-03}

\begin{exhibitbox}[三种财务口径不可拼接]
\small
\begin{tabularx}{\exhibitboxwidth}{L{3.1cm}R{2.5cm}R{2.5cm}L{3.6cm}X}
\toprule
\textbf{口径} & \textbf{2024收入} & \textbf{2024归母} & \textbf{可用场景} & \textbf{禁止事项} \\
\midrule
原600150法定旧口径 & 785.84 & 36.14 & 历史归一化与交易前公司观察 & 不作2025同比基数 \\
交易报告备考口径 & 另有备考数 & 交易分析口径 & 评估交易结构与范围 & 不替代法定实际 \\
2025年报重述比较期 & 1,333.51 & 42.20 & 2025同比及持续经营序列 & 不与旧股本EPS混用 \\
2025A持续经营实际 & 1,519.78 & 78.48 & 模型基期 & 不将范围扩大写成内生翻倍 \\
\bottomrule
\end{tabularx}
\sourcenote{CF-02、CF-03、CF-12。金额单位：亿元。}
\end{exhibitbox}

重组的真正价值只有在经营指标中出现才值得溢价：交付效率、费用率、毛利、少数股东归属和营运资金周转改善。规模跳升本身不创造每股价值，因为新增资产同时对应新增股份；这也是本报告不单独给“央企整合”估值倍数的原因。

\begin{riskbox}[口径失效风险]
若后续报告再次把legacy、交易备考和法定重述数据拼接，或用75.2562亿股重算2025A法定EPS，所有增长率、PE与目标价都会被污染。发现口径错配时应停止局部修补，回到重述利润表、加权分母和合并后期末股本整体重跑。
\end{riskbox}

因此，重组对估值的直接结论是“扩大了一个可审计的合并平台”，不是“自动获得协同溢价”。上市边界已经足以支持合并估值；协同是否值钱，留待交付、盈利和现金章节验证。
